\section*{Capítulo \ref{ch:g4:volcanoes-earthquakes-first-look} --- Volcanes y terremotos: una primera mirada}

\begin{solution}{exo:g4:volcanoes-earthquakes-first-look:1}
(a) \omterm{def:g4:volcanoes-earthquakes-first-look:eruption}{Lava}. (b) \omterm{def:g4:volcanoes-earthquakes-first-look:eruption}{Ceniza volcánica}. (c) Gas.
\end{solution}

\begin{solution}{exo:g4:volcanoes-earthquakes-first-look:2}
Un \omterm{def:g4:volcanoes-earthquakes-first-look:earthquake}{terremoto} es una sacudida repentina del suelo, causada por \omterm{def:g3:stones-pebbles-sand:rock}{rocas} que se
rompen y se deslizan en las profundidades. Dura segundos (un minuto o dos
en los más grandes), no horas.
\end{solution}

\begin{solution}{exo:g4:volcanoes-earthquakes-first-look:3}
Agáchate, cúbrete, agárrate.
\end{solution}

\begin{solution}{exo:g4:volcanoes-earthquakes-first-look:4}
Alrededor del \omterm{def:g2:water-on-earth:ocean}{océano} Pacífico.
\end{solution}

\begin{solution}{exo:g4:volcanoes-earthquakes-first-look:5}
(a) Verdadero: la mayoría solo la detectan los instrumentos. (b) Falso: la
\omterm{def:g4:volcanoes-earthquakes-first-look:eruption}{lava} se enfría y se endurece hasta convertirse en \omterm{def:g3:stones-pebbles-sand:rock}{roca} nueva. (c) Falso:
correr mientras el suelo tiembla es peligroso (puedes caerte, y de los
edificios caen cristales y objetos); agáchate, cúbrete y agárrate.
\end{solution}

\begin{solution}{exo:g4:volcanoes-earthquakes-first-look:6}
Viene de muy por debajo del \omterm{def:g4:volcanoes-earthquakes-first-look:volcano}{volcán}, donde la \omterm{def:g3:stones-pebbles-sand:rock}{roca} se ha fundido. Sube por el
centro del \omterm{def:g4:volcanoes-earthquakes-first-look:volcano}{volcán} y sale por arriba.
\end{solution}

\begin{solution}{exo:g4:volcanoes-earthquakes-first-look:7}
(a) Chile (también valen Perú, Ecuador o México): sus \omterm{def:g4:volcanoes-earthquakes-first-look:volcano}{volcanes} recorren todo
el país, cerca de la \omterm{def:g1:landscapes-around-us:water}{costa} del Pacífico. (b) Hawái, en medio del \omterm{def:g2:water-on-earth:ocean}{océano}
Pacífico, lejos del anillo sombreado.
\end{solution}

\begin{solution}{exo:g4:volcanoes-earthquakes-first-look:8}
$141 \div 12$ es un poco menos de 12 (porque $12 \times 12 = 144$): unos
12 grandes \omterm{def:g4:volcanoes-earthquakes-first-look:earthquake}{terremotos} al mes.
\end{solution}

\begin{solution}{exo:g4:volcanoes-earthquakes-first-look:9}
De 20, la cuarta parte es 5, así que tres cuartos son $3 \times 5 = 15$:
unos 15 en el anillo, y unos $20 - 15 = 5$ fuera de él.
\end{solution}

\begin{solution}{exo:g4:volcanoes-earthquakes-first-look:10}
Debería haberse agachado apoyando las manos y las rodillas, haberse
cubierto la cabeza y el cuello debajo de su mesa y haberse agarrado a ella
hasta que dejara de temblar. El marco de una puerta no le protege de los
objetos que caen o salen despedidos, y correr mientras el suelo tiembla
puede hacerle caer.
\end{solution}

\begin{solution}{exo:g4:volcanoes-earthquakes-first-look:11}
Fuera del anillo: $1214 - 837 = 377$ \omterm{def:g4:volcanoes-earthquakes-first-look:volcano}{volcanes}. Un tercio de 1\,200 es 400,
así que dos tercios de 1\,200 son 800. El anillo contiene 837, un poco más
de 800 de unos 1\,200: ``unos 2 de cada 3'' es una buena descripción.
\end{solution}
